\documentclass[11pt]{amsart}

\usepackage[T1]{fontenc}
\usepackage{lmodern}
\usepackage[margin=1.12in]{geometry}
\usepackage{amsmath,amssymb,mathtools}
\usepackage{booktabs}
\usepackage{tabularx}
\usepackage{enumitem}
\usepackage{xcolor}
\usepackage[colorlinks=true,linkcolor=blue!55!black,citecolor=blue!55!black,
  urlcolor=blue!55!black]{hyperref}

\newtheorem{theorem}{Theorem}[section]
\newtheorem{proposition}[theorem]{Proposition}
\newtheorem{lemma}[theorem]{Lemma}
\newtheorem{corollary}[theorem]{Corollary}
\theoremstyle{definition}
\newtheorem{definition}[theorem]{Definition}
\theoremstyle{remark}
\newtheorem{remark}[theorem]{Remark}

\newcommand{\CH}{\operatorname{CH}}
\newcommand{\divg}{\operatorname{div}}
\newcommand{\supp}{\operatorname{supp}}
\newcommand{\E}{\mathbb E}
\newcommand{\C}{\mathcal C}
\newcommand{\B}{\mathcal B}
\newcommand{\one}{\mathbf 1}
\newcommand{\edge}[2]{\mathbf e_{#1,#2}}

\title[Nonnegative flows for Chow-polynomial moments]
  {Nonnegative flows for Chow-polynomial moments of matroids}
\author{Anonymous}
\date{August 28, 2026}

\begin{document}

\begin{abstract}
Cheng and Lin expressed central-moment inequalities for the standard Chow
polynomial of a matroid as linear inequalities in flag numbers.  We prove
their Conjecture~3.30 in every rank by exhibiting its variance flag form as
an explicit nonnegative combination of their normalized flag-count
inequalities.  We also prove the case $k=4$ of their Conjecture~3.14: if
$M$ is a loopless matroid of rank $d+1$, then
\[
 \E\left[(X_M-d/2)^4\right]\le
 \begin{cases}
  1/16,&d=1,\\
  1/3,&d=2,\\
  (d+2)(5d+8)/240,&d\ge3.
 \end{cases}
\]
The fourth-moment result is strengthened to an identity expressing its flag
form as a nonnegative combination of the same normalized inequalities.
Both proofs use one mechanism: flag types form a support subgraph of the
Boolean lattice, the desired coefficient vector is a divergence, and
finite-poset summation by parts converts a nonnegative flow into the required
linear combination.  The variance flow uses peel-two splits.  For the fourth
moment, a signed combination of three primitive peel flows is converted into
a nonnegative peel-width-three flow by explicit zero-divergence reroutings.
All identities and inequalities in the proof are symbolic and uniform in
the rank.
\end{abstract}

\maketitle

\section{Introduction}

Throughout, matroids are loopless, as in the global convention of
Cheng and Lin~\cite[p.~2]{ChengLin}.  Let $M$ have rank $d+1$.  Its
standard Chow polynomial is
\[
  \CH_M(x)=\sum_{p=0}^{d}h_p(M)x^p,
  \qquad h_p(M)=\dim_{\mathbb Q}A^p(M),
\]
where $A^p(M)$ is the degree-$p$ part of its Chow ring.  Following
\cite[Definition~3.1, p.~8]{ChengLin}, the associated random variable
$X_M$ is defined by
$\Pr(X_M=p)=h_p(M)/\CH_M(1)$.  We follow the conventions and notation of
Cheng and Lin~\cite{ChengLin}.  Their moment conjectures compare $X_M$ to
the corresponding random variable $X_d$ for the Boolean matroid on
$d+1$ elements.  The latter has the Eulerian distribution.

The standard Chow ring was introduced in this combinatorial form by
Feichtner and Yuzvinsky~\cite{FY}; its Hodge theory was established by
Adiprasito, Huh, and Katz~\cite{AHK}.  Chow-polynomial formulas and related
Hilbert--Poincar\'e series have been developed further, for example in
\cite{FMSV}.  Cheng and Lin's flag expansion turns the moment comparisons
into inequalities in the flag numbers of $M$.  Our contribution is an
explicit common certificate for two such inequalities.

The first result settles their Conjecture~3.30.

\begin{theorem}[Cheng--Lin Conjecture~3.30]
\label{thm:intro-k2}
For every $d\ge1$, the variance flag form of Cheng and Lin is an explicit
nonnegative linear combination of the normalized flag-count inequalities in
their Lemma~3.29.  The combination is supported on Boolean cover pairs.
\end{theorem}

Our second result proves the fourth-moment instance of Cheng and Lin's
Conjecture~3.14.

\begin{theorem}[Fourth-moment comparison]
\label{thm:intro-k4}
Let $M$ be a loopless matroid of rank $d+1$, with $d\ge1$.  Then
\begin{equation}
 \E\left[(X_M-d/2)^4\right]\le
 \begin{cases}
  \dfrac1{16},&d=1,\\[4pt]
  \dfrac13,&d=2,\\[4pt]
  \dfrac{(d+2)(5d+8)}{240},&d\ge3.
 \end{cases}
 \label{eq:intro-k4-bound}
\end{equation}
Equivalently, this is Conjecture~3.14 of \cite{ChengLin} for $k=4$.
\end{theorem}

The proof gives more than the numerical comparison.  If $N_J$ denotes
the number of flags of flats of rank type $J\subseteq[d]$, $U_J$ the
corresponding Boolean flag number, and $a_J^{(4)}$ the fourth-moment
coefficient derived in Section~\ref{sec:k4-coefficient}, then there are
explicit $f_d(J,J')\ge0$, nonzero only for covers $J\lessdot J'$, such
that
\begin{equation}
 \sum_{J\subseteq[d]}a_J^{(4)}N_J
 =\sum_{J\lessdot J'}f_d(J,J')
 \left(\frac{N_J}{U_J}-\frac{N_{J'}}{U_{J'}}\right)\le0.
 \label{eq:intro-stronger}
\end{equation}
The equality in \eqref{eq:intro-stronger} is a formal identity of linear
forms; only the final inequality uses matroid theory.

The common structure is as follows.  A flag type $J$ is encoded by a
composition of $d+2$.  The coefficient vanishes unless every block has
size at least two, giving a support-composition DAG.  We orient Boolean
covers upward, from a composition to a refinement.  A flow whose divergence
is the coefficient vector yields, by summation by parts, a linear
combination of normalized flag-count differences.  The $k=2$ divergence
has a direct peel-two flow.  At $k=4$, primitive peel-two, peel-three, and
peel-four flows first give a signed certificate.  Differences of two paths
with common endpoints then reroute every negative peel-four occurrence.
A collision estimate proves that the positive peel-two and peel-three
reserves cannot be exhausted.

No bounded computation is used in these arguments.  We make no assertion
about equality cases, optimality of the displayed certificates, higher even
moments, or the augmented Chow polynomial.

\section{Flag types and support compositions}
\label{sec:notation}

Write $[d]=\{1,\ldots,d\}$.  For
$J=\{j_1<\cdots<j_m\}\subseteq[d]$, put $j_0=0$, $r=m+1$, and
\begin{equation}
 n_i=j_i-j_{i-1}\quad(1\le i\le m),
 \qquad n_r=d+2-j_m.
 \label{eq:composition}
\end{equation}
For $J=\varnothing$, this means $c(J)=(d+2)$.  Thus
$c(J)=(n_1,\ldots,n_r)$ is a composition of
\[
 N:=d+2
\]
whose final part is at least two.  Conversely, $J$ is the set of proper
prefix sums of this composition.  We call $n_r$ the final block and all
earlier blocks nonfinal.

Let $N_J=N_J(M)$ be the number of chains of nonempty proper flats of
$M$ whose ranks form $J$, with $N_\varnothing=1$.  The Boolean flag
number is
\begin{equation}
 U_J=\frac{(d+1)!\,n_r}{\prod_{i=1}^{r}n_i!}.
 \label{eq:UJ}
\end{equation}
In particular $U_J>0$.  Cheng and Lin's Notation~3.25 and
Proposition~3.26~\cite[p.~20]{ChengLin} give
\begin{equation}
 \CH_M(x)=\sum_{J\subseteq[d]}N_JP_J(x),\qquad
 P_J(x)=x^{-1}\prod_{i=1}^{r}f_{n_i}(x),
 \qquad f_n(x)=x+\cdots+x^{n-1}.
 \label{eq:flag-expansion}
\end{equation}
The formula \eqref{eq:UJ} is recorded on p.~21 of \cite{ChengLin}.

We say that $c(J)$ is on the \emph{support} when $n_i\ge2$ for every
$i$.  Indeed,
\begin{equation}
 P_J(1)=\prod_i(n_i-1),
 \label{eq:PJ1}
\end{equation}
so $P_J=0$ and every coefficient used below vanishes if any $n_i=1$.
The support is denoted $\C_N$.  It consists of the compositions of $N$
with all parts at least two and with the inherited final-block convention.

A Boolean cover $J\lessdot J\cup\{j\}$ splits one block of $c(J)$.
If the prefix sum before block $i$ is
$\pi_i=n_1+\cdots+n_{i-1}$ and $p=j-\pi_i$, the block $n_i$ is
replaced by $(p,n_i-p)$.  The allowed offsets are
\[
 1\le p\le n_i-1\quad(i<r),\qquad
 1\le p\le n_r-2\quad(i=r).
\]
The stricter final bound is exactly the condition $j\le d$.  A cover
between support states therefore splits a block into two parts at least two.
We orient all such edges from the coarser composition to the finer one.

On the support, let $Y_i$ be independent and uniform on
$\{1,\ldots,n_i-1\}$.  Dividing \eqref{eq:flag-expansion} by
\eqref{eq:PJ1} identifies the probability generating function of a flag
summand with that of
\begin{equation}
 F_J=\sum_iY_i-1.
 \label{eq:FJ}
\end{equation}
Since $\E Y_i=n_i/2$, all nonzero flag summands have the common mean
\begin{equation}
 \E F_J=\frac12\sum_i n_i-1=\frac d2.
 \label{eq:common-mean}
\end{equation}
This common mean permits central moments of the normalized Chow polynomial
to be computed term by term.

The only sign theorem about matroids used below is Cheng and Lin's
Lemma~3.29~\cite[p.~22]{ChengLin}:
\begin{equation}
 J\subseteq J'\quad\Longrightarrow\quad
 \frac{N_J}{U_J}\le\frac{N_{J'}}{U_{J'}}.
 \label{eq:CL-monotonicity}
\end{equation}

\section{Finite-poset summation by parts}
\label{sec:poset}

The passage from a combinatorial flow to a matroid inequality is formal.

\begin{lemma}[Finite-poset flow lemma]
\label{lem:poset-flow}
Let $P$ be a finite poset, orient each cover $u\lessdot v$ from $u$
to $v$, and let $f(u,v)$ be a flow on these edges.  Define
\[
 \divg f(u)=\sum_{u\to v}f(u,v)-\sum_{v\to u}f(v,u).
\]
For every function $x:P\to\mathbb R$,
\begin{equation}
 \sum_{u\in P}\divg f(u)x(u)
 =\sum_{u\lessdot v}f(u,v)\bigl(x(u)-x(v)\bigr).
 \label{eq:summation-by-parts}
\end{equation}
Consequently, if $f\ge0$, $x$ is isotone, and
$\divg f=w$, then
\[
 \sum_{u\in P}w(u)x(u)\le0.
\]
\end{lemma}

\begin{proof}
Expand the left side of \eqref{eq:summation-by-parts}.  A fixed edge
$u\to v$ occurs once with coefficient $+x(u)$ in the outflow at $u$
and once with coefficient $-x(v)$ in the inflow at $v$.  This proves the
identity.  Under the stated sign assumptions, every summand on the right is
nonpositive.
\end{proof}

For the Boolean lattice, take $x_J=N_J/U_J$.  Thus any nonnegative flow
with divergence $w_J=U_Ja_J$ proves
\begin{equation}
 \sum_Ja_JN_J
 =\sum_{J\lessdot J'}f(J,J')
 \left(\frac{N_J}{U_J}-\frac{N_{J'}}{U_{J'}}\right)\le0,
 \label{eq:generic-flag-certificate}
\end{equation}
where the equality is formal and the inequality follows from
\eqref{eq:CL-monotonicity}.

\section{The peel-two flow and Conjecture~3.30}
\label{sec:k2}

We first derive the variance coefficient.  Put
\[
 v(n)=n(n-2),\qquad V(c)=\sum_i v(n_i).
\]
By \eqref{eq:FJ}, independence, and the variance of a discrete uniform
variable,
\begin{equation}
 \E\left[(F_J-d/2)^2\right]=\frac{V(c(J))}{12}.
 \label{eq:block-variance}
\end{equation}
The Boolean Eulerian variable has variance $N/12=(d+2)/12$.  Normalize
\eqref{eq:flag-expansion} at $x=1$, use the common mean, and subtract the
Boolean variance.  We obtain
\begin{align}
 &6\CH_M(1)\left(
   \E[(X_M-d/2)^2]-\frac{N}{12}\right)\notag\\
 &\qquad=\sum_{J\subseteq[d]}N_J
 \prod_i(n_i-1)\frac{V(c(J))-N}{2}
 =\sum_{J\subseteq[d]}a_J^{(2)}N_J,
 \label{eq:k2-flag-form}
\end{align}
where
\begin{equation}
 a_J^{(2)}=\prod_i(n_i-1)
 \left(\frac12\sum_i n_i(n_i-3)\right).
 \label{eq:a2}
\end{equation}
This is the coefficient in Cheng and Lin's Corollary~3.28 and
Conjecture~3.30~\cite[p.~22]{ChengLin}.

For a support composition $c=(n_1,\ldots,n_r)$, define
\begin{equation}
 G(c)=U_c\prod_i(n_i-1)
 =(N-1)!\,n_r\prod_i\frac{n_i-1}{n_i!},
 \qquad
 \rho(n)=\frac{n(n-3)}2.
 \label{eq:k2-G-rho}
\end{equation}
Then $G(c)>0$, and the desired divergence is
\begin{equation}
 w^{(2)}(c)=U_ca_c^{(2)}=G(c)\sum_i\rho(n_i).
 \label{eq:k2-target}
\end{equation}

We now define the flow.  For every nonfinal block $m\ge4$, select the
split
\[
 m\longmapsto(2,m-2).
\]
For a final block $m\ge4$, select both
\[
 m\longmapsto(2,m-2),\qquad m\longmapsto(m-2,2).
\]
Each selected move receives raw flow $G(\text{child})$.  When a final
block $4$ is split into $(2,2)$, the two labels denote the same Boolean
cover edge, and their flows are added.  All unselected edges receive zero.
Denote the resulting edge flow by $f^{(2)}_d$.

\begin{lemma}[Peel-two ratios]
\label{lem:k2-ratios}
If a block $m\ge4$ is peeled as above, then
\[
 \frac{G(\text{child})}{G(\text{parent})}=
 \begin{cases}
  \dfrac{m(m-3)}2,&\text{nonfinal},\\[5pt]
  \dfrac{(m-2)(m-3)}2,&\text{final-left},\\[5pt]
  m-3,&\text{final-right}.
 \end{cases}
\]
The last two ratios sum to $\rho(m)$.
\end{lemma}

\begin{proof}
All unaffected factors in \eqref{eq:k2-G-rho} cancel.  For a nonfinal
split, the factor $(m-1)/m!$ is replaced by
$(1/2)(m-3)/(m-2)!$, giving $m(m-3)/2$.  For a final-left split, the
final-block prefactor changes from $m$ to $m-2$, multiplying the first
ratio by $(m-2)/m$.  For a final-right split, it changes from $m$ to
$2$, giving $m-3$.  Their sum is
$(m-3)((m-2)/2+1)=m(m-3)/2$.
\end{proof}

\begin{lemma}[Peel-two incoming edges]
\label{lem:k2-incoming}
For every support composition $c$, the selected move occurrences with
child $c$ are in bijection with the blocks of $c$ equal to $2$, and
each occurrence carries flow $G(c)$.
\end{lemma}

\begin{proof}
A nonfinal $2$-block is reversed by merging it with its right neighbor;
the resulting block has size at least four and its selected nonfinal-left
peel returns $c$.  A penultimate $2$-block is reversed in the same way
as a final-left peel, and a final $2$-block is reversed by merging it with
its left neighbor as a final-right peel.  These inverse operations are
unique.  If the final two blocks are both $2$, the two occurrences have
the same parent edge, but that edge was assigned the sum of both labeled
flows.  Hence multiplicity is preserved.  No selected edge meets an
off-support state.
\end{proof}

\begin{proposition}[Variance divergence]
\label{prop:k2-divergence}
For every $d\ge1$, $f^{(2)}_d\ge0$ and
$\divg f^{(2)}_d=w^{(2)}$ on the full Boolean lattice.
\end{proposition}

\begin{proof}
At a support composition $c$, Lemma~\ref{lem:k2-ratios} gives
\[
 \operatorname{outflow}(c)
 =G(c)\sum_{n_i\ge4}\rho(n_i).
\]
Lemma~\ref{lem:k2-incoming} gives
$\operatorname{inflow}(c)=G(c)\#\{i:n_i=2\}$.  Since
$\rho(2)=-1$ and $\rho(3)=0$, their difference is exactly
$G(c)\sum_i\rho(n_i)=w^{(2)}(c)$.  At an off-support state, both the
target and every incident selected flow vanish.  Nonnegativity follows from
$G>0$.  For $d=1$, the sole support state is $(3)$, and both sides are
zero.
\end{proof}

Combining Proposition~\ref{prop:k2-divergence} with
Lemma~\ref{lem:poset-flow} proves the identity
\[
 \sum_Ja_J^{(2)}N_J
 =\sum_{J\lessdot J'}f^{(2)}_d(J,J')
 \left(\frac{N_J}{U_J}-\frac{N_{J'}}{U_{J'}}\right).
\]
All coefficients are nonnegative.  This proves
Theorem~\ref{thm:intro-k2}, including the precise linear combination claim
of Cheng and Lin's Conjecture~3.30.

\section{Derivation of the fourth-moment coefficient}
\label{sec:k4-coefficient}

We next derive the coefficient to which the fourth-moment flow will be
matched.  Continue to write
\[
 v_i=n_i(n_i-2),\qquad V=\sum_i v_i,
 \qquad Q_2=\sum_i v_i^2.
\]
If $Z_i=Y_i-n_i/2$, direct finite summation gives
\begin{equation}
 \E Z_i=0,\qquad
 \E Z_i^2=\frac{v_i}{12},\qquad
 \E Z_i^4=\frac{v_i(3n_i^2-6n_i-4)}{240}.
 \label{eq:block-moments}
\end{equation}
The $Z_i$ are independent and centered.  Therefore
\begin{align}
 720\E[(F_J-d/2)^4]
 &=3\sum_i v_i(3n_i^2-6n_i-4)
    +30\sum_{i<j}v_iv_j\notag\\
 &=3\sum_i v_i(3v_i-4)+30\sum_{i<j}v_iv_j\notag\\
 &=9Q_2-12V+15(V^2-Q_2)\notag\\
 &=15V^2-6Q_2-12V.
 \label{eq:k4-block-expansion}
\end{align}
Thus the polynomial occurring below is a consequence of the flag expansion
and the moments of a finite uniform variable, not a definition chosen to
fit the flow.  Write $\mu_4(Z)=\E[(Z-d/2)^4]$ for the fourth central moment
of any variable in the present rank.

Let
\[
 \beta_d=\E[(X_d-d/2)^4],\qquad B_d=720\beta_d.
\]
Cheng and Lin's Corollary~3.13~\cite[pp.~13--14]{ChengLin}, together with
the two small Eulerian polynomials, gives
\begin{equation}
 B_1=45,\qquad B_2=240,\qquad
 B_d=3(d+2)(5d+8)=15N^2-6N\quad(d\ge3).
 \label{eq:Bd}
\end{equation}
Indeed, $1+x$ gives $\beta_1=1/16$, while
$1+4x+x^2$ gives $\beta_2=1/3$.

Normalize \eqref{eq:flag-expansion}, subtract $\beta_d$, and use
\eqref{eq:k4-block-expansion}.  Line by line,
\begin{align}
 &720\CH_M(1)\left(
   \E[(X_M-d/2)^4]-\beta_d\right)\notag\\
 &\quad=\sum_{J\subseteq[d]}N_JP_J(1)
 \left(720\E[(F_J-d/2)^4]-B_d\right)\notag\\
 &\quad=\sum_{J\subseteq[d]}N_J\prod_i(n_i-1)
 \left(15V_J^2-6Q_{2,J}-12V_J-B_d\right)\notag\\
 &\quad=\sum_{J\subseteq[d]}a_J^{(4)}N_J,
 \label{eq:k4-flag-form}
\end{align}
where
\begin{equation}
 a_J^{(4)}
 =720P_J(1)\bigl(\mu_4(F_J)-\mu_4(X_d)\bigr)
 =\prod_i(n_i-1)
 \left(15V_J^2-6Q_{2,J}-12V_J-B_d\right).
 \label{eq:a4}
\end{equation}
Consequently, with $G$ as in \eqref{eq:k2-G-rho},
\begin{equation}
 w^{(4)}(c):=U_ca_c^{(4)}
 =G(c)S_4(c),\qquad
 S_4(c)=15V(c)^2-6Q_2(c)-12V(c)-B_d.
 \label{eq:k4-target}
\end{equation}

\section{Primitive fourth-moment flows}
\label{sec:k4-primitives}

For $s\in\{2,3,4\}$, define a move multiset $F_s$ on the support DAG.
From every nonfinal block $m\ge s+2$, select
$m\mapsto(s,m-s)$.  From every final block $m\ge s+2$, select both
$m\mapsto(s,m-s)$ and $m\mapsto(m-s,s)$.  Give every selected move raw
flow $G(\text{child})$.  If the two final labels yield the same balanced
edge, their contributions are added.

Define
\begin{align}
 \tau_2(n)&=\frac{n(n-3)}2,&
 \tau_3(n)&=\frac{n(n-2)(n-4)}3,\notag\\
 \tau_4(n)&=\frac{n(n-2)(n-3)(n-5)}8,&
 T_s(c)&=\sum_i\tau_s(n_i).
 \label{eq:taus}
\end{align}

\begin{lemma}[Primitive divergences]
\label{lem:primitive-divergences}
For $d\ge3$ and $s\in\{2,3,4\}$,
\begin{equation}
 \frac{\divg F_s(c)}{G(c)}=T_s(c).
 \label{eq:primitive-divergence}
\end{equation}
\end{lemma}

\begin{proof}
For a nonfinal split $m\mapsto(s,m-s)$, factorial cancellation gives
\[
 \frac{G(\text{child})}{G(\text{parent})}
 =\binom ms\frac{(s-1)(m-s-1)}{m-1}.
\]
For a final block, the two orientations acquire the respective final-factor
ratios $(m-s)/m$ and $s/m$, whose sum is one.  Hence the sum of the two
final ratios equals the same expression.  Substitution of $s=2,3,4$
gives respectively $\tau_2(m),\tau_3(m),\tau_4(m)$.

Every occurrence of a block equal to $s$ has one incoming selected move:
merge it with the right neighbor when it is nonfinal and with the left
neighbor when it is final.  A balanced final pair is counted with the two
labels already built into $F_s$.  Each such incoming occurrence has raw
flow $G(c)$.  Now $\tau_s(s)=-1$, and $\tau_s$ vanishes at every other
support size below $s+2$.  Outflow minus inflow is therefore
$G(c)T_s(c)$.  The only sole-block exceptions are $c=(3)$ for $s=3$
and $c=(4)$ for $s=4$, which occur in $d=1,2$ and are treated
separately in Section~\ref{sec:small-d}.
\end{proof}

Let $F_{2,V}$ have the same selected occurrences as $F_2$, but give an
edge occurrence raw flow $G(\text{child})V(\text{child})$.

\begin{lemma}[Discrete product rule]
\label{lem:product-rule}
For $d\ge3$,
\begin{equation}
 \frac{\divg F_{2,V}}G
 =\frac{V(V-N)}2-6T_3-2V.
 \label{eq:product-rule}
\end{equation}
\end{lemma}

\begin{proof}
First, $T_2=(V-N)/2$.  A selected peel-two split
$m\mapsto(2,m-2)$ changes $V$ by $-4(m-2)$.  Moving the child
potential through the divergence in \eqref{eq:primitive-divergence} gives
\[
 \frac{\divg F_{2,V}}G
 =VT_2+\sum_i[-4(n_i-2)\tau_2(n_i)].
\]
The identity
\[
 -4(n-2)\tau_2(n)=-6\tau_3(n)-2n(n-2)
\]
now proves \eqref{eq:product-rule}.
\end{proof}

The further one-variable identity
\begin{equation}
 v(n)^2-9v(n)=8\tau_4(n)+18\tau_3(n)
 \label{eq:one-variable-k4}
\end{equation}
implies
\begin{equation}
 Q_2-9V=8T_4+18T_3.
 \label{eq:Q2-T}
\end{equation}

\begin{proposition}[Signed fourth-moment flow]
\label{prop:signed-flow}
For $d\ge3$, the signed flow
\begin{equation}
 f_0=30F_{2,V}+(30N-12)F_2+72F_3-48F_4
 \label{eq:signed-flow}
\end{equation}
satisfies $\divg f_0=w^{(4)}$.  Equivalently, a selected $F_2$
occurrence has coefficient $30V(\text{child})+30N-12$, an $F_3$
occurrence has coefficient $72$, and an $F_4$ occurrence has coefficient
$-48$, each multiplying $G(\text{child})$.
\end{proposition}

\begin{proof}
Using Lemmas~\ref{lem:primitive-divergences} and
\ref{lem:product-rule}, the normalized divergence of \eqref{eq:signed-flow}
is
\begin{align*}
 &30\left(\frac{V(V-N)}2-6T_3-2V\right)
 +(30N-12)\frac{V-N}{2}+72T_3-48T_4\\
 &\qquad=15V^2-108T_3-48T_4-66V-15N^2+6N.
\end{align*}
Equation \eqref{eq:Q2-T} changes this to
$15V^2-6Q_2-12V-15N^2+6N$, which equals $S_4$ by
\eqref{eq:Bd} for $d\ge3$.
\end{proof}

\section{Zero-divergence rerouting}
\label{sec:rerouting}

Only the $F_4$ occurrences in $f_0$ have negative coefficients.  We
eliminate them without changing divergence.  For an oriented edge $A\to B$
write $\edge AB$ for its unit flow.  The difference of the edge chains of
two directed paths with the same initial and terminal vertices has zero
divergence.

We use a dot to concatenate blocks, and allow $P,S$ to be empty support
compositions.  Consider first a nonfinal or final-left $F_4$ occurrence.
For $t\ge2$, set
\begin{equation}
 \begin{aligned}
 A&=P.(t+4).S,&B&=P.(4,t).S,\\
 C&=P.(2,t+2).S,&D&=P.(2,2,t).S.
 \end{aligned}
 \label{eq:left-square-vertices}
\end{equation}
The paths $A\to B\to D$ and $A\to C\to D$ have the same endpoints.
With $R=48G(B)$, add
\begin{equation}
 R\bigl(\edge AB+\edge BD-\edge AC-\edge CD\bigr).
 \label{eq:left-square}
\end{equation}
This cancels the negative $F_4$ occurrence on $A\to B$, adds flow on
$B\to D$, and subtracts only from selected $F_2$ occurrences.  Relative
to the normalizing factor $G(\text{child})$ of the consumed edge, the
subtractions are
\begin{equation}
 \begin{array}{c|c}
 \text{edge}&\text{normalized subtraction}\\ \hline
 A\to C&12(t+2)(t-1)\quad\text{if the original block is nonfinal},\\
 A\to C&12t(t-1)\quad\text{if the original block is final},\\
 C\to D&24.
 \end{array}
 \label{eq:left-square-costs}
\end{equation}
These identities follow directly by canceling the common factorial factors
in $48G(B)/G(C)$ and $48G(B)/G(D)$.

For a final-right $F_4$ occurrence, set
\begin{equation}
 A=P.(t+4),\quad B=P.(t,4),\quad
 C=P.(t+2,2),\quad D=P.(t,2,2).
 \label{eq:right-square-vertices}
\end{equation}
The analogous path-square correction cancels $A\to B$, adds flow on the
final $4\mapsto(2,2)$ edge $B\to D$, and consumes
$24(t+2)(t-1)$ normalized units from the final-right $F_2$ edge
$A\to C$.  It leaves the negative raw amount $R=48G(B)$ on the right
peel $C\to D$, namely $(t+2)\mapsto(t,2)$.

We transfer that remaining debt recursively.  At a stage with
$y=t-2j\ge4$, let $2^j$ denote $j$ consecutive blocks equal to $2$,
and put
\begin{equation}
 \begin{aligned}
 X_j&=P.2^j.(y+2,2),&Y_j&=P.2^j.(y,2,2),\\
 Z_j&=P.2^{j+1}.(y,2),&W_j&=P.2^{j+1}.(y-2,2,2).
 \end{aligned}
 \label{eq:cascade-vertices}
\end{equation}
Add
\begin{equation}
 R\bigl(\edge{X_j}{Y_j}+\edge{Y_j}{W_j}
       -\edge{X_j}{Z_j}-\edge{Z_j}{W_j}\bigr).
 \label{eq:cascade-square}
\end{equation}
This cancels the current negative right peel, adds flow on
$Y_j\to W_j$, consumes the selected left $F_2$ occurrence
$X_j\to Z_j$, and moves the same debt to $Z_j\to W_j$.  Each step lowers
$y$ by two.  If the process reaches $y=2$, the debt lands on the selected
$F_2$ edge $4\mapsto(2,2)$; if it reaches $y=3$, it lands on the
selected $F_3$ edge $5\mapsto(3,2)$.  Thus the process terminates.

Apply \eqref{eq:left-square} to every labeled nonfinal and final-left
$F_4$ occurrence, and apply the final-right correction and cascade to every
labeled final-right occurrence.  Balanced final moves retain both labels
until all corrections are made; in particular, the two occurrences of
$8\mapsto(4,4)$ are canceled separately.  Define $f_d$ to be the
result of adding all these corrections to $f_0$.  This is an explicit
finite prescription on every support DAG.  All vertices and edges appearing
in the path squares are support vertices and legal upward covers.  Every
correction has zero divergence, and therefore
\begin{equation}
 \divg f_d=w^{(4)}.
 \label{eq:corrected-divergence}
\end{equation}
All surviving flow is supported on peel-two and peel-three cover edges.

\section{Collision accounting and uniform nonnegativity}
\label{sec:nonnegative}

We prove that the simultaneous corrections never overdraw a positive edge.
The initial normalized reserve of a selected $F_2$ occurrence is
\begin{equation}
 H=30V(\text{child})+30N-12,
 \label{eq:F2-reserve}
\end{equation}
and that of a selected $F_3$ occurrence is $72$.

\begin{lemma}[Cascade collision bound]
\label{lem:collision}
Fix a selected nonfinal-left $F_2$ edge consumed by a cascade, and write
its affected block as $y+2$, with $y\ge2$.  Its possible cascade origins
are indexed by suffixes of the consecutive run of $2$-blocks immediately
preceding that block.  Their normalized contributions are $48r_j$, where
\begin{equation}
 r_j=2^j\frac{\rho_0(y+2j)}{\rho_0(y)},\qquad
 \rho_0(n)=\frac{n-1}{n!}.
 \label{eq:rj}
\end{equation}
They satisfy
\begin{equation}
 \frac{r_{j+1}}{r_j}
 =\frac{2}{(y+2j+2)(y+2j-1)},
 \qquad
 \sum_j r_j\le\frac{25}{16}.
 \label{eq:rj-bound}
\end{equation}
For a terminal $F_3$ edge, $y=3$ and
\begin{equation}
 \sum_j r_j\le\frac{79}{65}<\frac54.
 \label{eq:rj-F3}
\end{equation}
\end{lemma}

\begin{proof}
An origin is recovered uniquely by deleting a suffix of the indicated
finite $2$-run and replacing $y$ by $y+2j$.  Cancellation of all
unaffected factors in $G$ gives \eqref{eq:rj}.  Since
\[
 \frac{\rho_0(n+2)}{\rho_0(n)}
 =\frac1{(n+2)(n-1)},
\]
the ratio formula follows.  For $y\ge2$, the first ratio is at most
$1/2$ and each later ratio is at most $1/9$.  Hence
$\sum_j r_j\le1+(1/2)/(1-1/9)=25/16$.  At $y=3$, the corresponding
bounds are $1/5$ and $1/14$, giving
$1+(1/5)/(1-1/14)=79/65$.
\end{proof}

Thus all cascades together consume at most $75$ normalized units from a
fixed selected $F_2$ edge and strictly fewer than $60$ from a terminal
$F_3$ edge.  For completeness, we record why this is the entire collision
list.  Reverse a tagged correction from a consumed edge.  A top square
recovers its parent and affected block uniquely.  A bottom square first
removes the immediately preceding $2$-block.  A cascade removes a unique
suffix of the preceding $2$-run.  These alternatives are disjoint and
exhaustive.  Distinct affected blocks cannot define the same Boolean cover.
This occurrence-level classification also covers coincidences among the
$F_2,F_3,F_4$ labels at final block sizes $5,6,7$, and the doubled
balanced labels at sizes $4,6,8$: negative $F_4$ labels are canceled one
at a time, while coincident positive reserves are added.

\begin{proposition}[Uniform nonnegativity]
\label{prop:nonnegative}
For every $d\ge3$, every edge value of $f_d$ is nonnegative.
\end{proposition}

\begin{proof}
Consider first a selected nonfinal-left $F_2$ edge whose affected block
has size $M$.  For $M\ge6$, its unique possible top subtraction is at
most $12(M-2)(M-5)$, its possible bottom subtraction is $24$, and all
cascades together cost at most $75$.  The child has
$V\ge(M-2)(M-4)$, and the presence of a later final block gives
$N\ge M+2$.  After all subtractions, \eqref{eq:F2-reserve} is therefore
at least
\begin{align*}
 &30(M-2)(M-4)+30(M+2)-12\\
 &\qquad{}-12(M-2)(M-5)-24-75
 =18M^2-66M+69>0.
\end{align*}
For $M=4,5$, the initial reserves are at least $168,288$, respectively,
whereas the remaining possible consumption is at most $24+75=99$.

For a final-left $F_2$ edge of size $M\ge6$, no cascade consumes the
edge.  Total consumption is at most
$12(M-4)(M-5)+24$, while its reserve is at least
$30(M-2)(M-4)+30M-12$.  The remaining margin is
\[
 18M^2-42M-36>0.
\]
The cases $M=4,5$ are immediate from the same exact costs.

For a final-right $F_2$ edge of size $M\ge6$, its only subtraction is
the initial right square, of normalized size $24(M-2)(M-5)$.  The reserve
minus this cost is at least
\[
 6M^2+18M-12>0.
\]
Again $M=4,5$ are immediate.  At final $M=4$, the left and right labels
coincide, so the positive reserves add.

Finally, a terminal $F_3$ edge starts with reserve $72$ and loses
strictly less than $60$ by Lemma~\ref{lem:collision}.  Every correction
not listed above adds rather than subtracts flow.  The three margin
polynomials are already positive at $M=6$, where their values are
$321,360,312$, and their first differences are positive thereafter.
Hence every aggregated edge flow is nonnegative.
\end{proof}

\section{Boundary dimensions and the matroid theorem}
\label{sec:small-d}

For $d=1$, the support has the single composition $(3)$.  Here
$S_4(3)=0$, so $w^{(4)}=0$ and we set $f_1=0$.

For $d=2$, the support DAG consists of
\[
 (4)\longrightarrow(2,2).
\]
Directly from the definitions,
\[
 G(4)=G(2,2)=3,\qquad S_4(4)=240,\qquad S_4(2,2)=-240.
\]
Assigning raw flow $720$ to this edge gives the required divergence.

Together with Propositions~\ref{prop:signed-flow} and
\ref{prop:nonnegative}, these cases prove the all-dimensional combinatorial
statement.

\begin{theorem}[Uniform fourth-moment flow]
\label{thm:uniform-flow}
For every $d\ge1$, the prescription above defines a nonnegative flow
$f_d$ on the support-composition DAG such that
\[
 \divg f_d(c)=w^{(4)}(c)
\]
at every support state.  Every used edge splits a block as
$n\mapsto(p,n-p)$ with $\min\{p,n-p\}\le3$.
\end{theorem}

Extend $f_d$ by zero to the full Boolean lattice.  Off the support,
$w^{(4)}=0$, and no selected edge is incident, so the divergence equality
continues to hold everywhere.  Apply Lemma~\ref{lem:poset-flow} with
$x_J=N_J/U_J$.  We obtain the formal identity
\begin{equation}
 \sum_Ja_J^{(4)}N_J
 =\sum_{J\lessdot J'}f_d(J,J')
 \left(\frac{N_J}{U_J}-\frac{N_{J'}}{U_{J'}}\right).
 \label{eq:final-flow-identity}
\end{equation}
Cheng and Lin's monotonicity \eqref{eq:CL-monotonicity} makes every bracket
nonpositive.  Since $f_d\ge0$, the right side is at most zero.  Equation
\eqref{eq:k4-flag-form} and $\CH_M(1)>0$ then yield
$\E[(X_M-d/2)^4]\le\beta_d$.  Substitution of \eqref{eq:Bd} proves
Theorem~\ref{thm:intro-k4} and the stronger identity
\eqref{eq:intro-stronger}.

This argument applies to every loopless matroid in the standard Chow
setting of \cite{ChengLin}.  No realizability, field, or simplicity
hypothesis is introduced, and no additional matroid-theoretic reduction is
needed.

\section{The common support-DAG mechanism}
\label{sec:common-mechanism}

The two proofs separate into three layers:
\begin{enumerate}[label=\textup{(\roman*)}]
 \item Cheng and Lin's flag expansion converts a moment difference into a
 coefficient vector $a_J$, supported where every composition block is at
 least two.
 \item A moment-specific, purely combinatorial construction writes
 $w_J=U_Ja_J$ as the divergence of a nonnegative upward flow.
 \item The finite-poset identity converts that divergence into a sum of
 normalized flag-count differences, and Cheng and Lin's Lemma~3.29 supplies
 their sign for matroids.
\end{enumerate}
The support reduction, cover dictionary, and final implication are identical
for $k=2$ and $k=4$.  What changes is only the middle layer.  For the
variance, the block polynomial $\rho(n)$ is itself the outgoing
peel-two ratio, and incoming flow records the negative value
$\rho(2)=-1$.  For the fourth moment, the $V^2$ and $Q_2$ interactions
require the primitive product rule and the identity
\eqref{eq:one-variable-k4}; zero-divergence path squares then turn the
resulting signed flow into a nonnegative one.

\section{Dependencies and machine verification}
\label{sec:dependencies}

Table~\ref{tab:dependencies} separates external matroid input from the new
combinatorial arguments.

\begin{table}[ht]
\caption{Logical dependencies.}
\label{tab:dependencies}
\small
\begin{tabularx}{\textwidth}{@{}
  >{\raggedright\arraybackslash}p{.28\textwidth}
  >{\raggedright\arraybackslash}X
  >{\raggedright\arraybackslash}p{.23\textwidth}@{}}
\toprule
Input & Role & Status \\
\midrule
Cheng--Lin Definition~3.1 and global convention
 & Defines the standard Chow polynomial and the loopless setting
 & External notation \\
Cheng--Lin Notation~3.25 and Proposition~3.26
 & Gives \eqref{eq:composition} and the flag expansion
   \eqref{eq:flag-expansion}
 & External matroid theorem \\
Boolean flag formula $U_J$
 & Identifies the normalized target $w_J=U_Ja_J$
 & External formula from \cite[p.~21]{ChengLin} \\
Cheng--Lin Corollary~3.13
 & Gives the Boolean fourth moment for $d\ge3$
 & External; $d=1,2$ computed directly \\
Finite uniform moments
 & Derive the $k=2$ and $k=4$ coefficients
 & Elementary algebra proved here \\
Peel flows, path-square corrections, and reserve bounds
 & Construct nonnegative flows with the required divergences
 & New combinatorial results proved here \\
Finite-poset flow lemma
 & Converts divergence into the exact linear-form identities
 & Formal incidence identity proved here \\
Cheng--Lin Lemma~3.29
 & Supplies the final sign $N_J/U_J\le N_{J'}/U_{J'}$
 & Sole external sign input \\
Additional matroid reduction
 & None required
 & Not applicable \\
\bottomrule
\end{tabularx}
\end{table}

Exact symbolic tests accompany the proof.  They check the finite-uniform
moment expansions, the coefficient bridge, all primitive divergence and
product identities, the occurrence-level path-square corrections, collision
accounting, reserve inequalities, boundary cases, and the final finite-poset
identity.  These tests use integer or rational arithmetic.  Finite support
DAG evaluations are regression checks only; the proof of
Theorem~\ref{thm:uniform-flow} is the symbolic argument in
Sections~\ref{sec:k4-primitives}--\ref{sec:nonnegative}.

The $k=2$ combinatorial core has also been formalized in Lean~4 with
mathlib: the definitions of $G,w$, the selected edge flow, nonnegativity,
the full divergence identity, and the abstract summation-by-parts conclusion
are machine checked.  Cheng and Lin's Lemma~3.29 is supplied there as an
explicit monotonicity hypothesis; no matroid or Chow-ring theory is
formalized.  The axiom report for the final theorem contains only
\texttt{propext}, \texttt{Classical.choice}, and \texttt{Quot.sound}, the
standard mathlib dependencies.  The $k=4$ proof has not been formalized in
Lean.  Neither the exact test suite nor the Lean development replaces any
paper proof above.

\begin{thebibliography}{99}

\bibitem{ChengLin}
Ronnie Cheng and Wangyang Lin,
\emph{Inequalities for Chow polynomials and Chern numbers of matroids},
arXiv:2603.21680v3 [math.CO], 29 April 2026,
\url{https://arxiv.org/abs/2603.21680v3}.

\bibitem{FY}
E. M. Feichtner and S. Yuzvinsky,
\emph{Chow rings of toric varieties defined by atomic lattices},
Invent. Math. \textbf{155} (2004), no.~3, 515--536,
\href{https://doi.org/10.1007/s00222-003-0327-2}{doi:10.1007/s00222-003-0327-2}.

\bibitem{AHK}
K. Adiprasito, J. Huh, and E. Katz,
\emph{Hodge theory for combinatorial geometries},
Ann. of Math. (2) \textbf{188} (2018), no.~2, 381--452,
\href{https://doi.org/10.4007/annals.2018.188.2.1}{doi:10.4007/annals.2018.188.2.1}.

\bibitem{FMSV}
L. Ferroni, J. P. Matherne, M. Stevens, and L. Vecchi,
\emph{Hilbert--Poincar\'e series of matroid Chow rings and intersection
cohomology},
Adv. Math. \textbf{449} (2024), Paper No.~109733,
\href{https://doi.org/10.1016/j.aim.2024.109733}{doi:10.1016/j.aim.2024.109733}.

\end{thebibliography}

\end{document}
